%%
%%   CV following the TENK Researcher's Curriculum Vitae Template (2020),
%%   implemented with the official LaTeX template of the Research Council
%%   of Finland: https://www.aka.fi/globalassets/10rahoitus/ansioluettelot/cv-template_latex.tex
%%   THE STYLE AND HEADERS HAVE BEEN PRESET ACCORDING TO THE RULES.
%%
\documentclass[a4paper,12pt]{article}

% Language support (including Nordic letters)
\usepackage[T1]{fontenc}
\usepackage[english]{babel}

% Math support
\usepackage[intlimits]{amsmath}
\usepackage{amsthm}

% Symbol support
\usepackage{amssymb}
\usepackage{eurosym}

% Enabling pdf
\usepackage[pdftex]{graphicx}
\usepackage{xcolor}

% Item list extras for reduced spacing between items
\usepackage{enumitem}
\setlist{nosep}

% Hyperlink support (only for external hyperlinks)
\usepackage[implicit=false]{hyperref}
\hypersetup{
    colorlinks=true,
    urlcolor=blue,
}

%-----------------------------------------------------
% ACADEMY SETTINGS: DO NOT ALTER THESE SETTINGS!
%-----------------------------------------------------

% ONLY CALIBRI/CARLITO ARE ALLOWED AS FONT, FONT SIZE 12 PT DEFINED ABOVE. The font is used to ensure equal settings between LaTeX and Word templates.
% Find the font package here: https://www.ctan.org/tex-archive/fonts/carlito
\usepackage[sfdefault]{carlito}

% LINE SPACING 1 SUPPORT AND SET SPACE
\usepackage{setspace}
\singlespacing

% PAGE MARGIN SUPPORT AND SET PAGE MARGINS 20 MM
\usepackage[left=20mm, top=20mm, bottom=20mm, right=20mm]{geometry}

% SETTING SPACING FOR SECTION, SUBSECTION AND PARAGRAPH
\usepackage{titlesec}
\titleformat*{\section}{\bfseries}
\titleformat*{\subsection}{\bfseries}
\titlespacing{\section}{0 cm}{6pt}{0pt}
\titlespacing{\subsection}{0 cm}{6pt}{0pt}
\setlength{\parskip}{6pt}

% ALIGNING TEXT TO LEFT
\raggedright

% SETS PAGE NUMBERS AND HEADER
\usepackage{fancyhdr}
\pagestyle{fancyplain}
\fancyfoot{}
\fancyhead{}
\renewcommand{\headrulewidth}{0pt}
\rhead{\fancyplain{}{\thepage}({\pageref{lastpage}}) }

%-----------------------------------------------------

\begin{document}

\section*{Personal details and date of CV}

\begin{itemize}
    \item Surname: Arora
    \item First names: Nikunj
    \item Researcher identifier: \href{https://orcid.org/0009-0000-7113-1196}{ORCID 0009-0000-7113-1196}
    \item Date of CV: 18.10.2026
\end{itemize}

\section*{Degrees}

\begin{itemize}
    \item 06/2021: Master of Science, major subject Computer Science (specialisation in Artificial Intelligence), University of Oulu, Oulu, Finland. Thesis: \emph{Augmenting Virtual Reality Telepresence Experience Using Self-Avatar}.
    \begin{itemize}
        \item Contact details of the organisation granting the degree: P.O. Box 8000, FI-90014 University of Oulu, Finland (\href{https://www.oulu.fi/en}{www.oulu.fi/en})
    \end{itemize}
    \item 06/2017: Bachelor of Technology, Electronics and Communication Engineering, Guru Gobind Singh Indraprastha University, New Delhi, India.
\end{itemize}

\section*{Language skills}

\begin{itemize}
    \item Hindi: native
    \item English: C2 (self-assessment, Europass scale)
    \item Finnish: A1 (self-assessment, Europass scale)
\end{itemize}

\section*{Current employment}

\begin{itemize}
    \item 06/2021--05/2026: Specialist Software Engineer, Qt Group, Oulu, Finland (full-time). Software development for the Qt framework and developer tooling: authored Claude-compatible LLM skills to guide developers through onboarding with AI coding agents; created and maintained the Qt for Android Studio plugin and the companion Qt Gradle plugin; Android/AOSP platform work for automotive IVI systems, including hidden API integration, activity embedding, and the Media Sessions API for Qt for Android Automotive; automated Android Automotive AOSP image builds via a build server.
\end{itemize}

\section*{Previous work experience}

\begin{itemize}
    \item 2020: Research Assistant, Center for Ubiquitous Computing, University of Oulu, Finland. Built a Unity crowd simulation for a telepresence robot to navigate through in a virtual world. The work later evolved into the M.Sc. thesis, which focused on VR telepresence and self-avatar embodiment through a 20-participant pilot study with psychometric measures (presence, co-presence, embodiment) and real-time behavioural measures; the pilot was subsequently extended into the pre-registered 62-participant study published in IEEE TVCG (see Research output).
    \item 2017--2019: Game Developer, Sameer Hazari Studios Pvt.\ Ltd., New Delhi, India (full-time). Built the full-stack AWS backend and frontend for a 3D commerce pipeline; automated 3D model conversion across formats (including USDZ for Apple devices); developed an iOS augmented reality app and Unreal Engine 4 virtual reality experiences for immersive 3D product previews; applied a deep learning style-transfer model to auto-generate product art variations; built a browser-side 3D rendering pipeline producing product videos client-side.
\end{itemize}

\newpage
\section*{Research output}

Publications (1):

\begin{itemize}
    \item N.~Arora, M.~Suomalainen, M.~Pouke, E.~G.~Center, K.~J.~Mimnaugh, A.~P.~Chambers, S.~Pouke, S.~M.~LaValle. \emph{Augmenting Immersive Telepresence Experience with a Virtual Body.} IEEE Transactions on Visualization and Computer Graphics, 28(5):2135--2145, 2022. \href{https://doi.org/10.1109/TVCG.2022.3150473}{DOI: 10.1109/TVCG.2022.3150473}
\end{itemize}

Methods, software, and tools developed:

\begin{itemize}
    \item Real-time 360\textdegree{} video streaming pipeline with self-avatar embodiment for an HMD telepresence experiment (Unity).
    \item Python scripts for automating quantitative data collection from the raw accelerometer and gyroscope signals of the HMD and controllers during the telepresence experiment.
\end{itemize}

\section*{Scientific and societal impact}

\begin{itemize}
    \item Open science: preprint of the peer-reviewed telepresence study made openly available (\href{https://arxiv.org/abs/2202.00900}{arXiv:2202.00900}).
\end{itemize}

\section*{Other}

\begin{itemize}
    \item Open-source projects (\href{https://github.com/su-nikunj}{github.com/su-nikunj}): SyncEngine, a Kotlin CLI simulating two-way synchronisation between a central server and 100\,000 simulated devices over a REST API using concurrency; Stereo-Image-ZNCC, a C++ stereo correspondence implementation; an OpenGL renderer with a debugging inspector for real-time graphics experiments.
    \item Other expertise: programming in Python, C\#, C++, Kotlin, Java, JavaScript, and TypeScript; development with Unity and Unreal Engine 4; machine learning; AWS; AI-assisted software development tooling. Interest in human--AI interaction, including trust and privacy in AI-mediated interaction and the human factors of immersive and AI-driven systems.
\end{itemize}

%% Page counter for header
\label{lastpage}

\end{document}
